% =============================================================================
% Anhang A: Konformitätsmatrix
% =============================================================================

\chapter{Konformitätsmatrix}
\label{chap:anhang-konformitaet}

Die folgende Tabelle enthält die am Primärtext geprüfte Einzelzuordnung aller 67 Zeilen der
Konformitätsmatrix. Die Summen und ihre Auswertung stehen in \secref{sec:soll-ist};
\tabref{tab:eval-sollist-summe} fasst sie zusammen.

Die Tabelle trennt Anwendbarkeit (Spalte A), Implementierungsstand (Spalte I) und technische
Erfüllbarkeit (Spalte T). In Spalte A steht J für anwendbar und N für nicht anwendbar. In Spalte I
steht V für vollständig und P für teilweise umgesetzt; nicht anwendbare Zeilen erhalten N. In Spalte T
steht V für vollständig, P für teilweise und X für technisch nicht erfüllbar; \texttt{n.a.}
kennzeichnet Zeilen ohne technisches Urteil. Die Belegspalte kürzt die Module ab: W steht für
\file{src/wire}, O für \file{src/options}, R für \file{src/recv}, F für \file{src/frag}, A für die
Schnittstelle unter \file{src/api} und K für die Kampagnenartefakte; jede Modulangabe schließt die
zugehörigen Tests ein. Code und eigene Tests gelten dabei nicht als voneinander unabhängige
Belegfamilien.\footnote{Die Einzelzuordnung liegt als \file{thesis/daten/konformitaet"=kategorien.csv}
im Repositorium der Arbeit; ein Prüfskript erzeugt Tabelle und Datei aus demselben Bestand, in dem
die drei Urteile je Zeile getrennt gespeichert sind, und prüft die Summen.}

% Generiert von thesis/tikz/gen/ff1-kategorien.py aus dem dortigen Datenbestand.
% Nicht von Hand aendern; Aenderungen im Skript oder in der Datenschicht vornehmen.
\begingroup
\footnotesize
\setlength{\tabcolsep}{2pt}
\renewcommand{\arraystretch}{1.00}
\begin{longtable}{@{}r l L{4.45cm} L{2.00cm} c c c L{4.20cm}@{}}
\caption{Prüfung der 67 Matrixzeilen gegen RFC 9868 und den Implementierungsstand}
\label{tab:eval-sollist}\\
\toprule
\textbf{Z.} & \textbf{Sec.} & \textbf{Normtext, gekürzt} & \textbf{Stufe} & \textbf{A} &
\textbf{I} & \textbf{T} & \textbf{Beleg} \\
\midrule
\endfirsthead
\toprule
\textbf{Z.} & \textbf{Sec.} & \textbf{Normtext, gekürzt} & \textbf{Stufe} & \textbf{A} &
\textbf{I} & \textbf{T} & \textbf{Beleg} \\
\midrule
\endhead
\bottomrule
\endfoot
162 & 10 & Ungültige UDP-Länge verwerfen und melden & MUST & J & V & V & R; FR-49 \\
163 & 7 & Surplus Area hinter dem Ende gemäß UDP-Länge & Festlegung & J & V & V & W; FR-04 \\
164 & 8 & Ganze Surplus Area nutzen; OCS an erster Zweiergrenze & Festlegung & J & V & V & W, O; FR-04/19 \\
165 & 8 & Pad vor OCS null, sonst Optionen verwerfen & MUST & J & V & V & W, R; FR-05; S-04 \\
166 & 9 & OCS über Bereich und 16-Bit-Längensummand & Festlegung & J & V & V & O; FR-19; Wire-Prüfung \\
167 & 9 & OCS nicht null, wenn UDP-Prüfsumme nicht null & MUST & J & V & V & O; FR-21; S-06 \\
168 & 9 & Nichtnull-OCS als Voreinstellung & MUST & J & V & V & O, A; FR-21 \\
169 & 9 & Bei OCS-Fehler Optionen und Surplus verwerfen & MUST & J & V & V & R; FR-20; S-05 \\
170 & 9/14 & Gültige Nutzdaten trotz OCS-Fehler standardmäßig liefern & MUST & J & V & V & R; FR-38; S-06 \\
171 & 10 & TLV-Rahmung und erweitertes Längenformat & Festlegung & J & V & V & O; FR-07/09 \\
172 & 10 & NOP und EOL ohne Längenfeld & Festlegung & J & V & V & O; FR-08 \\
173 & 10 & Länge unter Formatminimum verwirft alle Optionen & MUST & J & V & V & O, R; FR-10; S-11a \\
174 & 10 & Unterlauf oder Überlauf macht den Bereich ungültig & MUST & J & V & V & O, R; FR-10/50; S-11b \\
175 & 10 & Bekannte Option unter Mindestlänge verwirft alle Optionen & MUST & J & V & V & O, R; FR-11/12 \\
176 & 10 & Ab 255 erweitern; kleinste Kodierung bevorzugen & MUST\slash SHOULD & J & V & V & O; FR-09/14 \\
177 & 10/14 & Optionen in Drahtreihenfolge verarbeiten & MUST & J & V & V & O, R; FR-13 \\
178 & 10 & Acht Pflichtoptionen erkennen und bei Konfiguration erzeugen & MUST & J & V & V & O; FR-06/14/22 bis 31 \\
179 & 10 & Unbekannte SAFE-Optionen still übergehen & MUST & J & V & V & O, R; FR-11/12; S-10 \\
180 & 10 & Defektes FRAG wie nicht unterstützte UNSAFE-Option & MUST & J & V & V & R, F; FR-27 \\
181 & 10 & Bestimmte Optionen nur einmal; erste Instanz gilt & SHOULD\slash MUST & J & V & V & O, R; FR-15 \\
182 & 10 & NOP darf sich zur Ausrichtung wiederholen & MAY & J & V & V & O; FR-16 \\
183 & 10 & Mehrere FRAG machen den Optionsbereich ungültig & MUST & J & V & V & R, F; FR-29 \\
184 & 10 & Bei UNSAFE Nutzdaten leer und Inhalt im FRAG & MUST & J & V & V & R, F; FR-28/39 \\
185 & 10 & Erste unbekannte UNSAFE beendet die Verarbeitung & MUST & J & V & V & R; FR-39; S-20 \\
186 & 10 & Pflichtoptionen außer NOP/EOL vor anderen SAFE senden & MUST\slash MAY & J & V & V & O, A; FR-14 \\
187 & 11.1 & Bei Restplatz EOL zuletzt; danach null senden & MUST & J & V & V & O; FR-17; S-08 \\
188 & 11.1 & Optionale Nullprüfung nach EOL samt Verwerfregel & MAY\slash MUST & N & N & n.a. & FR-18; nicht gewählt \\
189 & 11.2 & Höchstens sieben aufeinanderfolgende NOP; dauerhafte Läufe melden & SHOULD & J & V & V &
  O, R; FR-16; NFR-05 \\
190 & 11.3 & APC als CRC32C; Fehlschlag angezeigt zustellen & Festlegung\slash SHOULD & J & V & V & O, R; FR-22/23 \\
191 & 11.3 & Unbekannte APC-Länge wie fehlgeschlagene APC & MUST & J & V & V & O, R; FR-12/23 \\
192 & 11.4 & FRAG-Länge 10 oder 12 abhängig von RDOS & Festlegung & J & V & V & F; FR-26; S-19 \\
193 & 11.4 & Bei FRAG sind UDP-Nutzdaten leer & MUST & J & V & V & R, F; FR-28 \\
194 & 11.4 & Mindestens zwei Fragmente, jedes in 1500 Byte, reassemblieren & MUST & J & V & V & F; FR-35; S-frag \\
195 & 11.4 & Kennung über den Timeout eindeutig; Erzeugung empfohlen & MUST\slash SHOULD & J & V & V & F; FR-31; S-46 \\
196 & 11.4 & Überlappung abbrechen und ohne ICMP verwerfen & MUST & J & V & V & F; FR-32/48; S-45 \\
197 & 11.4 & Exakte Duplikate dürfen verworfen werden & MAY & J & V & V & F; FR-32; S-45 Z0 \\
198 & 11.4 & Standardtimeout höchstens zwei Minuten; kein ICMP & SHOULD\slash MUST & J & V & V & F; FR-33; S-42 \\
199 & 11.4 & Speicher begrenzen und nicht socketübergreifend teilen & SHOULD & J & P & V &
  F, A; FR-34; NFR-06; eigene Korrektur von f847895 \\
200 & 11.4 & Einzelne Fragmente nie an den Nutzer geben & MUST & J & V & V & R, F; FR-31 \\
201 & 11.4 & Fehlgeschlagene Reassemblierung ohne Leerdatagramm & SHOULD & J & V & V & R, F; FR-33; S-49 \\
202 & 11.5 & MDS verarbeiten; MDS begrenzt das Senden nicht & MUST & J & V & V & O, A; FR-24 \\
203 & 11.6 & MRDS-Mindestwerte und Voreinstellungen für IPv4 & MUST & J & V & V & O, F; FR-35 \\
204 & 11.7 & REQ/RES ohne Autoantwort; RES-Token aus empfangenem REQ & Festlegung\slash MUST & J & V & V &
  O, A; FR-25; Vertrag \\
205 & 11.7 & Automatische REQ/RES-Schicht standardmäßig aus & MUST & J & V & V & A; keine solche Schicht \\
206 & 11.8 & TIME optional; verwendete Zeitwerte nie null & MAY\slash MUST & N & N & n.a. &
  nicht gewählt; S-23 Parser \\
207 & 11.9 & AUTH ist reserviert & reserviert & N & N & n.a. & nicht definiert \\
208 & 11.10 & EXP bei Wahl mit Mindestlänge vier & MUST & N & N & n.a. & nicht gewählt \\
209 & 12 & UCMP und UENC reserviert; UEXP als UNSAFE-Experimentoption & reserviert & N & N & n.a. & nicht definiert \\
210 & 12 & UNSAFE nur in Fragmenten; unbekannte Art verwirft Daten & MUST & J & V & V & R, F; FR-39 \\
211 & 13 & Entwurfsregeln für neue Optionsarten & MUST & N & N & n.a. & keine neue Optionsart \\
212 & 14 & Empfangsfolge Prüfsumme, OCS, Optionen, Zustellung & Festlegung & J & V & V & R; FR-36 \\
213 & 14 & Vorhandene Pflichtoptionen verarbeiten & MUST & J & V & V & R; FR-06/36 \\
214 & 14 & Andere Optionsarten dürfen ignoriert werden & MAY & J & V & V & R, A; gewählt \\
215 & 14 & Nutzdaten bei SAFE unabhängig vom Ergebnis liefern & MUST & J & V & V & R; FR-38 \\
216 & 14 & FRAG/NOP/EOL intern; übrige Status bereitstellen & MUST & J & P & V &
  R, A; FR-37; Parameter nicht immer verfügbar \\
217 & 14 & Optionsüberlauf nach Abschnitt 10 behandeln & MUST & J & V & V & O, R; FR-50; Erratum 8834 \\
218 & 15 & Empfangs-API: je Paket und Fragment gefordert oder aus & Normform & J & P & V & A; FR-44 \\
219 & 15 & Fehlende Pflichtoption still verwerfen und melden & MUST\slash SHOULD & J & P & V & A; FR-44 \\
220 & 15 & Schalter für alle Optionsdatagramme; Standard verarbeiten & Normform\slash MUST & J & V & V & A; FR-44 \\
221 & 15 & Optionen und Status für Nutzer verfügbar & Normform\slash MUST & J & P & V &
  R, A; FR-37; Parameter nicht immer verfügbar \\
222 & 15 & Sende-API: Auswahl, Geltung, Mindestlänge, Maximalfragment & Normform & J & P & V & A; SendConfig/Options \\
223 & 15/25 & Keine direkte Ordnungs- oder Fragmentgrenzensteuerung & Leitlinie & N & N & n.a. & A; kanonisiert \\
224 & 16/19 & Transit unverändert (nicht FF1); SAFE-Fehler ignorieren und Nutzdaten liefern & MUST & J & V & V &
  R, A; FR-38 für Endpunkt; Transit: K, nicht als I bewertet \\
225 & 25 & Bei DoS-Bedenken Grenzen für Optionen, NOP und Fragmente & bed. SHOULD & J & P & V & R, F; NFR-04/05/06/08 \\
226 & 25.2 & Bei Kanalbedenken Optionen in unabhängiger Referenzordnung & bed. SHOULD & N & N & n.a. &
  nicht gewählt; kanonische Sendereihenfolge genügt nicht \\
227 & 23 & Multicast und Broadcast gesondert betrachten & Leitlinie & N & N & n.a. & nur Unicast \\
228 & 26 & Namensregeln für neue SAFE- und UNSAFE-Arten & MUST & N & N & n.a. & keine neue Art \\
\midrule
\multicolumn{8}{@{}l@{}}{Implementierung: 50 V, 7 P, 10 N. Technisch: 57 V, 0 P, 0 nicht erfüllbar; 10 n.a.} \\
\end{longtable}
\endgroup

